% =====================================================================
% ch03c —— 第 3 章 §3.2 续（原书 p.103）＋ §3.3 群 J(X)（p.104–p.109）
%           ＋ 第 3 章章末评注
% =====================================================================

% [原文 p.103]
% 承接 p.102 末句：ここで、$u$ が $\tilde K(S^{2n})$ を生成するとし、
\begin{proof}[证明（续）]
\begin{equation*}
  f^*(u)=\sum_{i=n}^{n+k}a_ix^i\text{。}
\end{equation*}
$f$ 是收缩，故 $a_n=1$。由 $\psi^k(u)=k^nu$ 知 $f^*(u)$ 必为 $y^n$ 的倍数，于是
\begin{equation*}
  \sum_{i=n}^{n+k}a_ix^i=\lambda y^n
\end{equation*}
成立。限制到 $S^{2n}$ 上便知 $\lambda=1$。所以在
\begin{equation*}
  y^n=(\log(1+x))^n
\end{equation*}
中，从 $x^n$ 到 $x^{n+k}$ 的所有系数都如所求为整数。
% [Tutor Note: 限制到 $S^{2n}$ 时 $\tilde K(S^{2n})$ 中 $x^{n+1}=0$，故 $i^*(x^i)=0$（$i>n$）；同理 $i^*(y)$ 与 $i^*(x)$ 在 $K(S^{2n})$ 中相等，这正是 $\lambda=1$ 的来源。]
\end{proof}

\begin{remark}{}
  定理 3.2.10 也给出收缩存在的充分条件，这一点由 Adams 与 Grant Walker 证明（Math. Proc. Camb. Phil. Soc. 61 (1) (1965) 81--103）。
\end{remark}

再设给定映射 $f: S^{2m+2n-1}\to S^{2m}$。此时可用下述方式给 $f$ 配一个不变量 $e(f)\in\Q/\Z$。

取 $X$ 为 $f$ 的映射锥，$i: S^{2m}\to X$ 为包含映射，$j: X\to S^{2m+2n}$ 为把 $S^{2m}$ 压成一点的映射。设 $u$ 生成 $\tilde K^0(S^{2m+2n})$，$v$ 生成 $\tilde K^0(S^{2m})$，并设 $x\in\tilde K^0(X)$ 满足 $i^*(x)=v$。令 $y=j^*(u)$。此时对任意 $k$，
\begin{equation*}
  \psi^k(x)=k^mx+a_ky
\end{equation*}
成立。前面已经证明 $\psi^k\psi^\ell=\psi^\ell\psi^k$，因此
\begin{equation*}
  k^n(k^m-1)a_\ell=\ell^n(\ell^m-1)a_k
\end{equation*}
成立。于是
\begin{equation*}
  e(f)=\frac{a_k}{k^n(k^m-1)}\in\Q
\end{equation*}
% [原文 p.104]
上式在固定 $x$ 之后是良定义的。$x$ 若只改变 $y$ 的整数倍，$e(f)$ 就只改变一个整数，因此 $e(f)\in\Q/\Z$ 良定义。$e:\pi_{2n+2m-1}(S^{2m})\to\Q/\Z$ 是群同态，这一点作为基本练习题留给读者。这个 $e(f)$ 是非常强的不变量。

\section{群 $J(X)$}

\takeaway{本节引入向量丛的纤维同伦等价与稳定化，得到群 $J(X)$；再用 Adams 算子 $\psi^k$ 造出 $\sigma^k:K(X)\to G_k$，把 $J$ 的核圈在子群 $H_k$ 中。}

本节中为简单起见，假定 $X$ 连通。先引入向量丛之间的一种等价关系，叫做纤维同伦等价（fibre homotopy equivalent）。这个等价关系在同伦论中颇受关注。设 $E$, $E'$ 是空间 $X$ 上的向量丛，且 $E$ 与 $E'$ 上都给了 Hermite 度量。此时称 $E$ 与 $E'$ 纤维同伦等价，是指存在映射 $f: S(E)\to S(E')$ 与 $g: S(E')\to S(E)$，二者与到 $X$ 的投影可交换，并且 $gf$、$fg$ 作为保纤维映射都与恒等映射同伦。这显然在 $X$ 上向量丛的同构类集合上定义了一个等价关系。
% [Tutor Note: 定义只用到球面丛 $S(E)$ 与投影 $S(E)\to X$，与向量丛的加法结构无关；因此 $E$ 与 $E\oplus V$ 一般并不等价，必须再做稳定化才得到群。]

纤维同伦等价是可加的。即，若 $E$, $E'$ 分别与 $F$, $F'$ 纤维同伦等价，则 $E\oplus E'$ 与 $F\oplus F'$ 纤维同伦等价。这由下述事实推出：$S(E\oplus E')$ 可以看成两个纤维空间 $S(E)$, $S(E')$ 的纤维联接（fibre join）。一般地，$\pi: Y\to X$ 与 $\pi': Y'\to X$ 的纤维联接，定义为在满足 $t\in I$、$\pi(y)=\pi'(y')$ 的三元组 $(y,t,y')$ 上引入等价关系
\begin{equation*}
  (y,0,y_1')\sim(y,0,y_2'),\qquad (y_1,1,y')\sim(y_2,1,y')
\end{equation*}
所得的空间。

称两个向量丛 $E$, $E'$ 稳定纤维同伦等价，是指存在平凡丛 $V$, $V'$，使 $E\oplus V$ 与 $E'\oplus V'$ 纤维同伦等价。$X$ 上所有稳定纤维同伦等价类组成的集合构成半群，记作 $J(X)$。每个向量丛 $E$ 都有补丛 $F$，使
% [原文 p.105]
$E\oplus F$ 成为平凡丛，故 $J(X)$ 是群。于是，映射
\begin{equation*}
  \Vect(X)\longrightarrow J(X)
\end{equation*}
扩张为满同态
\begin{equation*}
  K(X)\longrightarrow J(X)
\end{equation*}
这个映射也记作 $J$。

设 $E$, $E'$ 是向量丛，其球面丛的投影分别为 $\pi: S(E)\to X$, $\pi': S(E')\to X$。此时，Thom 复形 $X^E$, $X^{E'}$ 分别是映射 $\pi$, $\pi'$ 的映射锥本身。因此，若 $E$ 与 $E'$ 纤维同伦等价，则 $X^E$ 与 $X^{E'}$ 具有相同的同伦型。另一方面，若 $E$ 是 $n$ 维平凡丛，则 $X^E=S^{2n}(X^+)$。因此，要证明 $J(E)\neq0$，只需证明 $X^E$ 与 $X^+$ 的悬垂不具有相同的稳定同伦型。
% [Tutor Note: 稳定同伦型在悬垂下不变，而 $S^{2n}(X^+)$ 与 $X^+$ 稳定同伦等价，所以这里可以用 $X^+$ 代替平凡丛的 Thom 复形。]

下面说明，为给出 $J(E)=0$ 的必要条件，如何使用 §3.2 的算子 $\psi^k$。由 Thom 同构（推论 2.7.12）可知，$\tilde K(X^E)$ 是由 $\lambda_E$ 生成的自由 $K(X)$ 模。于是，对每个 $k$，存在唯一的元素 $\rho^k(E)\in K(X)$ 使
\begin{equation*}
  \psi^k(\lambda_E)=\lambda_E\rho^k(E)
\end{equation*}
成立。把 §3.2 证明的基本类 $\lambda_E$ 的乘性与 $\psi^k$ 保积这一事实合起来，便得
\begin{equation*}
  \rho^k(E\oplus E')=\rho^k(E)\cdot\rho^k(E')\text{。}
\end{equation*}
又，取 $E=1$，并设 $\beta$ 为周期性定理的同构映射，回想
\begin{equation*}
  \psi^k\circ\beta=k\beta\circ\psi^k
\end{equation*}
便知
\begin{equation*}
  \rho^k(1)=k\text{。}
\end{equation*}
% [Tutor Note: 这里 $\lambda_E\in\tilde K(X^E)$ 是 Thom 类；$\tilde K(X^E)$ 是秩 1 的自由 $K(X)$ 模，所以 $\psi^k(\lambda_E)$ 必是 $\lambda_E$ 的 $K(X)$-倍数，$\rho^k(E)$ 由此唯一确定。]
% [原文 p.106]
这里，令 $\Q_k=\Z[1/k]$ 为所有分母是 $k$ 的幂的分数组成的 $\Q$ 的子环。此时，令
\begin{equation*}
  \sigma^k(E)=k^{-n}\rho^k(E)\qquad(n=\dim E)
\end{equation*}
便得到同态
\begin{equation*}
  \sigma^k: K(X)\longrightarrow G_k
\end{equation*}
其中 $G_k$ 是 $K(X)\otimes\Q_k$ 的单位组成的乘法群。这里，若 $E$ 与平凡丛纤维同伦等价，则存在 $u\in\tilde K(X^E)$ 使 $\psi^k(u)=k^nu$。令 $u=\lambda_Ea$，可知
\begin{equation*}
  \psi^k(\lambda_E)\cdot\psi^k(a)=k^n\lambda_Ea
\end{equation*}
于是
\begin{equation*}
  \sigma^k(E)\cdot\psi^k(a)=a
\end{equation*}
成立。进一步，限制到一点上，$a$ 在增广映射下的像是 1，故 $a$ 与 $\psi^k(a)$ 都是 $G_k$ 的单位。于是可写
\begin{equation*}
  \sigma^k(E)=\frac{a}{\psi^k(a)}\in G_k
\end{equation*}
由于 $\sigma^k(E)$ 只依赖 $E$ 的稳定类，下述命题得证。

\begin{proposition*}{3.3.1}
  设 $H_k\subset G_k$ 是由所有形如 $a/\psi^k(a)$（$a$ 取遍 $K(X)$ 的单位）的元素生成的子群。此时，
  \begin{equation*}
    \sigma^k: K(X)\longrightarrow G_k
  \end{equation*}
  把 $J$ 的核映入 $H_k$，因此诱导同态
  \begin{equation*}
    J(X)\longrightarrow G_k/H_k\text{。}
  \end{equation*}
\end{proposition*}

% [原文 p.107]
要用命题 3.3.1，需要能计算 $\sigma^k$，或者同样地，能计算 $\rho^k$。现在，
\begin{equation*}
  \rho^k\in\operatorname{Op}(K)
\end{equation*}
是算子。它在增广映射下的像已知，所以剩下只须确定它在直线丛的组合上的取值。由乘性，只须对直线丛 $L$ 确定 $\rho^k(L)$ 即可。

\begin{lemma*}{3.3.2}
  对直线丛 $L$，
  \begin{equation*}
    \rho^k[L]=\sum_{j=0}^{k-1}[L]^j
  \end{equation*}
  成立。
\end{lemma*}

\begin{proof}
由命题 2.7.1 与 2.7.2，$\tilde K(X^L)$ 可以描述为 $K(P(L\oplus1))$ 的 $K(X)$ 子模，并由 $n=1-[L][H]$ 生成。$K(P(L\oplus1))$ 的结构当然由主定理 2.2.1 给出。于是，
\begin{align*}
  \psi^k(u) &= 1-[L^k][H^k] \\
  &= (1-[L][H])\left\{\sum_{j=0}^{k-1}[L^j][H^j]\right\} \\
  &= u\sum_{j=0}^{k-1}[L^j]
\end{align*}
这里最后一步用了 $(1-[L][H])(1-[H])=0$。因此
\begin{equation*}
  \psi^k(\lambda_L)=\lambda_L\left\{\sum_{j=0}^{k-1}[L^j]\right\}
\end{equation*}
从而
\begin{equation*}
  \rho^k(L)=\sum_{j=0}^{k-1}[L^j]
\end{equation*}
得证。
\end{proof}

% [原文 p.108]
作为例子，考虑 $2n$ 维实射影空间 $X=P_{2n}(\R)$。如命题 2.7.7 所示，$\tilde K(X)$ 是以 $x=[L]-1$ 为生成元的 $2^n$ 阶循环群，其中 $L$ 是标准线丛。其乘法结构由关系 $[L]^2=1$ 看出（这来自 $L$ 附属于群 $\Z/2$）。取 $k=3$，则
\begin{equation*}
  \psi^3(x)=[L^3]-1=x
\end{equation*}
而命题 3.3.1 中定义的群 $H_3$ 成为平凡群。另一方面，用引理 3.3.2 可得
\begin{align*}
  \sigma^3(mx) &= \rho^3(mx)=(\rho^3(x))^m=(\rho^3[L])^m\cdot3^{-m} \\
  &= 3^{-m}(1+[L]+[L]^2)^m \\
  &= (1+x/3)^m \\
  &= 1+\sum_{i=1}^{m}(-1)^{i-1}\frac{2^{i-1}}{3^i}\binom{m}{i}x
    \qquad (x^2=-2x) \\
  &= 1+\frac12\left(1-\left(1-\frac23\right)^m\right)x \\
  &= 1+3^{-m}\left(\frac{3^m-1}{2}\right)x
\end{align*}
% [Tutor Note: 由 $[L]^2=1$ 得 $x^2=([L]-1)^2=2-2[L]=-2x$，这是把 $(1+x/3)^m$ 的所有高次项降为一次的根据。]
可知。因此，若 $J(mx)=0$，则 $3^m-1$ 必须被 $2^{n+1}$ 整除。这只有在 $2^{n-1}$ 整除 $m$ 时才会发生。因此
\begin{equation*}
  J:\tilde K(P_{2n}(\R))\longrightarrow J(P_{2n}(\R))
\end{equation*}
的核的阶至多为 2。
% [Tutor Note: $x$ 的阶为 $2^n$，故 $m$ 只须按模 $2^n$ 考虑；在 $\Z/2^n$ 中满足 $2^{n-1}\mid m$ 的只剩 $m\equiv0$ 与 $m\equiv2^{n-1}$，这就是“核的阶至多 2”的来源。]
实际上，这一结果可以用实 K 理论改进，改进后的结果是球面上向量场问题之解的基础。

实际上，定理 3.2.10 所考虑的问题，是现在所考虑的更一般问题的特殊情形。事实上，容易看出空间 $P_{n+k}(\C)/P_{n-1}(\C)$ 是 $P_k(\C)$ 上向量丛 $nH$ 的 Thom 空间。因此，定理 3.2.10 的结论可以解释为关于 $J[H]\in J(P_k(\C))$ 的阶的断言。定理 3.2.10 的证明手法与本节所用本质相同。关键在于，不把 $P_{n+k}(\C)/P_{n-1}(\C)$ 等当成各自的空间逐一处理，
% [原文 p.109]
而是作为 Thom 空间一次考虑，用统一的手法处理。

关于这一方向上 $J(X)$ 的更多细节，可以举出 J. F. Adams 的一系列论文 On the groups $J(X)$（\textit{Topology} 1963--）作为参考。

\chapterreview{
\textbf{全章地位。}第 3 章把第 2 章建立的 K 群从上同调论的框架推进到“运算”的层面。第 1、2 章回答 K 群是什么、怎样计算，本章则研究 $K(X)$ 上的自然变换本身，再用这些变换生产不变量。这一转向使 K 理论从关于结构的学问变成能算、能判别的工具，故第 3 章是全书由结构走向应用的分水岭：向前它依赖周期性定理与 Thom 同构，向后它直接连到附录 B 的幂算子。

\textbf{知识脉络。}§3.1 走 Grothendieck 的路线：先由 $\lambda_t$ 的乘性得到外幂运算 $\lambda^k$，再由 $\lambda_{-t}$ 与对称幂 $S_t$ 互逆的关系导出 $\gamma$ 算子与维数 $\dim_K$、$\dim_\gamma$，并用若干对称函数的引理证明 $\tilde K(X)$ 幂零。定理 3.1.7 把算子环 $\operatorname{Op}(K_1,K)$ 完全刻画为形式幂级数环，等于说算子就是对称函数的形式版本。§3.2 取对数导数，把乘性的 $\lambda$ 运算换成加性的 Adams 算子 $\psi^k$。§3.3 让 $\psi^k$ 作用在 Thom 复形的 Thom 类上，得到 $\rho^k$ 与 $\sigma^k$，把 $J(X)$ 的核夹进由 $a/\psi^k(a)$ 生成的子群 $H_k$。

\textbf{核心结论。}$\psi^k$ 是环同态，满足 $\psi^k\psi^\ell=\psi^{k\ell}$，还对素数 $p$ 有 $\psi^p(x)\equiv x^p$。在 $K(X)\otimes\Q$ 上它按 $k^i$ 作用，于是 $K(X)\otimes\Q$ 依“Betti 次数”分解为固有空间；命题 3.2.7 进一步把过滤商与 Betti 群之和等同起来，并给出 Euler 特征数 $\chi(X)=\dim_\Q(K^0(X)\otimes\Q)-\dim_\Q(K^1(X)\otimes\Q)$ 的表达式。定理 3.2.3 用这些同余定出 Hopf 不变量为奇值时只允许 $n=1,2,4$；定理 3.2.10 与不变量 $e(f)\in\Q/\Z$ 给出 $P_{n+k}(\C)/P_{n-1}(\C)$ 收缩到 $S^{2n}$ 的整性障碍。§3.3 把这些技术汇总成 $\ker J\subset H_k$，使 $J(X)$ 可以计算，$P_{2n}(\R)$ 一例算出核的阶至多 2。

\textbf{辨析与提醒。}务必分清 $\lambda^k$ 与 $\psi^k$：前者乘性，$\lambda_t(x+y)=\lambda_t(x)\lambda_t(y)$；后者加性且乘性，是环同态。$\psi_t$ 的定义中参数是 $-t$，又带因子 $-t\,d/dt$，符号错一处便得不到 $\psi^k([L])=[L]^k$。特征分解只在 $K(X)\otimes\Q$ 上成立，整环 $K(X)$ 上不能直接对角化，所以 §3.2 谈固有空间时处处先张量 $\Q$。所谓“Betti 群”由 $\psi^k$ 的固有空间造出，日译本 p.100 的译注已提醒它与通常上同调群的构造有别，只在有限性条件下才重合。$\sigma^k$ 只在 $E$ 的稳定类上有定义，取值落在乘法群 $G_k$ 而非 $K(X)$，因为要用到除法 $k^{-n}$。

\textbf{评注。}本章的技术核心是“用一个乘性运算的对数导数换取加性运算”，这个配方后来成为 Adams 型算子的一般框架，在 Adams 谱序列、广义上同调论与 $\lambda$ 环理论中反复出现。§3.3 的 $J(X)$ 只是第一个回报：Adams 此后关于 $J(X)$ 的论文沿同一路线把 $J(X)$ 的阶与 Bernoulli 数联系起来，最终通向 Adams 猜想与球面稳定同伦群的图像。就本书的编排而言，第 3 章的算子技术在附录 B 中被系统推广，而 $J(X)$ 与 Hopf 不变量的讨论则标出了本书在拓扑应用上的边界。
}

% CHUNK ch03c DONE
